% HADROS-CASCADE Pipeline Documentation
% Compile with: pdflatex hadros_cascade_pipeline.tex (twice for cross-references)

\documentclass[11pt,a4paper]{article}

\usepackage[utf8]{inputenc}
\usepackage[T1]{fontenc}
\usepackage[english]{babel}
\usepackage{amsmath,amssymb,amsfonts}
\usepackage{geometry}
\usepackage{graphicx}
\usepackage{xcolor}
\usepackage{booktabs}
\usepackage{array}
\usepackage{longtable}
\usepackage{hyperref}
\usepackage{enumitem}
\usepackage{fancyhdr}
\usepackage{titlesec}
\usepackage{microtype}
\usepackage{mdframed}
\usepackage{listings}
\usepackage{caption}
\usepackage{subcaption}
\usepackage{lmodern}

\geometry{
  top=2.5cm, bottom=2.5cm,
  left=2.5cm, right=2.5cm,
  headheight=14pt
}

% Colors
\definecolor{darkgreen}{RGB}{26,46,36}
\definecolor{brown}{RGB}{140,88,0}
\definecolor{codeblue}{RGB}{0,80,160}
\definecolor{codegray}{RGB}{100,100,100}
\definecolor{lightbg}{RGB}{245,247,243}
\definecolor{boxborder}{RGB}{180,200,172}

% Header/footer
\pagestyle{fancy}
\fancyhf{}
\fancyhead[L]{\small\textcolor{darkgreen}{\textbf{HADROS-CASCADE}}}
\fancyhead[R]{\small\textcolor{codegray}{Scientific Pipeline Documentation}}
\fancyfoot[C]{\small\thepage}
\renewcommand{\headrulewidth}{0.4pt}

% Section colors
\titleformat{\section}{\large\bfseries\color{darkgreen}}{{\thesection.}}{0.5em}{}[\vspace{-2pt}\rule{\linewidth}{0.6pt}]
\titleformat{\subsection}{\normalsize\bfseries\color{darkgreen!80!black}}{{\thesubsection}}{0.5em}{}
\titleformat{\subsubsection}{\normalsize\bfseries\color{brown}}{{\thesubsubsection}}{0.5em}{}

% Code listings
\lstset{
  basicstyle=\small\ttfamily,
  keywordstyle=\color{codeblue}\bfseries,
  commentstyle=\color{codegray}\itshape,
  backgroundcolor=\color{lightbg},
  frame=single,
  rulecolor=\color{boxborder},
  breaklines=true,
  columns=flexible,
}

% Physics box environment
\newmdenv[
  backgroundcolor=lightbg,
  linecolor=boxborder,
  linewidth=1pt,
  innerleftmargin=10pt,
  innerrightmargin=10pt,
  innertopmargin=8pt,
  innerbottommargin=8pt,
  skipabove=6pt,
  skipbelow=6pt,
]{physbox}

\hypersetup{
  colorlinks=true,
  linkcolor=darkgreen,
  citecolor=brown,
  urlcolor=codeblue,
  pdftitle={HADROS-CASCADE: Scientific Pipeline Documentation},
  pdfauthor={HADROS-CASCADE},
}

% ─────────────────────────────────────────────────────────────────────────────
\begin{document}

\begin{titlepage}
  \centering
  \vspace*{2cm}
  {\color{darkgreen}\rule{\linewidth}{2pt}}\\[0.8cm]
  {\Huge\bfseries\color{darkgreen} HADROS-CASCADE}\\[0.3cm]
  {\Large\color{brown} Scientific Pipeline Documentation}\\[0.5cm]
  {\color{darkgreen}\rule{\linewidth}{2pt}}\\[1.5cm]

  {\large UHE Neutrino Deep Inelastic Scattering\\
  in Kerr Black Hole Spacetime:\\[0.3cm]
  From Incoming Geodesic to Hadronic Cascade\\
  and Backward Camera Observation}\\[2cm]

  \begin{tabular}{ll}
    \textbf{Pipeline:} & HADROS-CASCADE \\[4pt]
    \textbf{Physics:}  & UHE $\nu$ DIS $+$ GEANT4 $+$ Kerr geodesics \\[4pt]
    \textbf{Backends:} & POWHEG-DIS $+$ PYTHIA~8 $+$ GEANT4 \\[4pt]
    \textbf{Date:}     & June 2026 \\
  \end{tabular}

  \vfill
  {\small\color{codegray}
  Gravitational Lensing $\cdot$ Kerr Null Geodesics $\cdot$ NLO DIS $\cdot$ Hadronic Showers\\
  ZAMO Tetrads $\cdot$ GBW/IIM Saturation Models $\cdot$ Backward Ray Tracing
  }
\end{titlepage}

\tableofcontents
\newpage

% ─────────────────────────────────────────────────────────────────────────────
\section{Overview and Physical Motivation}

HADROS-CASCADE is a Monte Carlo pipeline designed to simulate ultra-high-energy (UHE) neutrino interactions in the vicinity of a rotating (Kerr) black hole embedded in an accreting environment --- a collapsar. In HADROS-CASCADE, secondary particles are currently mapped with a particle-ray association camera: cascade products are associated with sampled Kerr rays by spatial and local angular criteria. This is not full transport of secondary particles to a distant observer.

The pipeline couples four major physics domains:
\begin{enumerate}[leftmargin=2cm]
  \item \textbf{General Relativity:} Kerr metric, null geodesic integration (backward ray tracing), ZAMO orthonormal tetrad frame.
  \item \textbf{High-Energy QCD:} NLO DIS cross sections (GBW/IIM saturation models), POWHEG-DIS event generation, PYTHIA~8 parton shower and hadronization.
  \item \textbf{Particle Transport:} Full hadronic and electromagnetic cascade via GEANT4 in a local flat-space box at the interaction point.
  \item \textbf{Particle-Ray Association Mapping:} Backward Kerr geodesic sampling, cascade-origin/particle-ray association, interaction-probability weighting.
\end{enumerate}

Figure~\ref{fig:pipeline_overview} summarizes the execution order.

\begin{figure}[h!]
\centering
\begin{mdframed}[backgroundcolor=lightbg, linecolor=boxborder]
\small\centering
\[
\underbrace{\text{Kerr Geodesics}}_{\text{Step 1}}
\;\longrightarrow\;
\underbrace{\text{Interaction Points}}_{\text{Step 2}}
\;\longrightarrow\;
\underbrace{\text{POWHEG+PYTHIA}}_{\text{Step 3}}
\;\longrightarrow\;
\underbrace{\text{GEANT4 Cascade}}_{\text{Step 4}}
\]
\[
\;\longrightarrow\;
\underbrace{\text{Backward Camera}}_{\text{Step 5}}
\;\longrightarrow\;
\underbrace{\text{Ray$\leftrightarrow$Event Link}}_{\text{Step 6}}
\;\longrightarrow\;
\underbrace{\text{GBW/IIM Weighting}}_{\text{Step 7}}
\;\longrightarrow\;
\underbrace{\text{Figures \& Dashboard}}_{\text{Step 8}}
\]
\end{mdframed}
\caption{Ordered execution of the HADROS-CASCADE pipeline. Each arrow represents a handoff of data (JSONL, binary, CSV) between steps.}
\label{fig:pipeline_overview}
\end{figure}

% ─────────────────────────────────────────────────────────────────────────────
\section{Step 1: Kerr Null Geodesic Tracing}
\label{sec:geodesics}

\subsection{Physical Role}

The first step establishes the photon/neutrino path geometry. A camera is placed at observer position $(r_\mathrm{obs},\,\theta_\mathrm{obs})$ in Boyer--Lindquist (BL) coordinates. For each pixel $(i,j)$ of the $N_x \times N_y$ detector, a null ray is fired \emph{backward} (from observer toward the source region) and integrated through the Kerr metric until it either crosses the outer torus boundary $r_\mathrm{max}$ or is captured by the black hole. The geodesic records the trajectory, cumulative column density, and optical depth at each step.

\subsection{The Kerr Metric}

In Boyer--Lindquist coordinates $(t, r, \theta, \phi)$, with $c = G = M_\bullet = 1$ (gravitational units, so distances in $r_g = GM_\bullet/c^2$), the Kerr line element is \cite{Kerr1963, Misner1973}:
\begin{equation}
  ds^2 = -\left(1 - \frac{2r}{\Sigma}\right)dt^2 - \frac{4ar\sin^2\theta}{\Sigma}\,dt\,d\phi
         + \frac{\Sigma}{\Delta}\,dr^2 + \Sigma\,d\theta^2
         + \frac{A\sin^2\theta}{\Sigma}\,d\phi^2,
  \label{eq:kerr}
\end{equation}
where the auxiliary functions are:
\begin{align}
  \Sigma &= r^2 + a^2\cos^2\theta, \label{eq:sigma}\\
  \Delta &= r^2 - 2r + a^2, \label{eq:delta}\\
  A &= (r^2 + a^2)^2 - a^2\Delta\sin^2\theta. \label{eq:A}
\end{align}

The spin parameter $a \in [0,\,1)$ (in units of $M_\bullet$) characterizes frame dragging; $a=0$ reduces to the Schwarzschild metric. The outer event horizon is located at:
\begin{equation}
  r_+ = 1 + \sqrt{1-a^2}.
\end{equation}

The non-zero contravariant metric components needed for geodesic integration are:
\begin{align}
  g^{tt} &= -\frac{A}{\Sigma\Delta}, \quad
  g^{t\phi} = -\frac{2ar}{\Sigma\Delta}, \quad
  g^{rr} = \frac{\Delta}{\Sigma}, \quad
  g^{\theta\theta} = \frac{1}{\Sigma}, \quad
  g^{\phi\phi} = \frac{\Sigma - 2r}{A\sin^2\theta}.
\end{align}

\subsection{Null Geodesic Equations}

Geodesics are integrated in the Hamiltonian formulation \cite{MTW, Fuerst2004}. Define the 8-dimensional phase-space vector $(x^\mu, p_\mu)$ where $p_\mu$ is the covariant four-momentum (with $p_0 < 0$ for future-directed rays). The geodesic equations are:
\begin{equation}
  \frac{dx^\mu}{d\lambda} = g^{\mu\nu} p_\nu,
  \qquad
  \frac{dp_\mu}{d\lambda} = -\frac{1}{2}\,\frac{\partial g^{\alpha\beta}}{\partial x^\mu}\,p_\alpha p_\beta,
  \label{eq:geodesic}
\end{equation}
subject to the null constraint:
\begin{equation}
  H = \frac{1}{2}\,g^{\mu\nu} p_\mu p_\nu = 0.
  \label{eq:null}
\end{equation}

Two constants of motion reduce the integration to 4 first-order ODEs. The conserved energy $E = -p_t$ and angular momentum $L = p_\phi$ along the symmetry axis give:
\begin{align}
  \frac{dt}{d\lambda} &= g^{tt}p_t + g^{t\phi}p_\phi = -g^{tt}E + g^{t\phi}L,\\
  \frac{d\phi}{d\lambda} &= g^{t\phi}p_t + g^{\phi\phi}p_\phi = g^{t\phi}E + g^{\phi\phi}L.
\end{align}
Carter's constant $\mathcal{Q}$ closes the system \cite{Carter1968}:
\begin{equation}
  \mathcal{Q} = p_\theta^2 + \cos^2\theta\left(a^2(m^2 - p_t^2) + \frac{p_\phi^2}{\sin^2\theta}\right),
\end{equation}
with $m = 0$ for null rays.

\subsection{Numerical Integration}

The code integrates Eq.~\eqref{eq:geodesic} using a fixed-step RK4 scheme (or adaptive RK5/Cash--Karp), with step size $h = \Delta r / (dr/d\lambda)$. At each step the integrator evaluates the metric and its partial derivatives via analytical formulae in \texttt{src/kerr\_metric.cpp}.

\subsection{Column Density Accumulation}

While tracing, the integrator accumulates the column density through the torus:
\begin{equation}
  \Lambda(s) = \int_0^s \rho\!\left(r(\lambda),\,\theta(\lambda)\right) \frac{dl}{d\lambda}\,d\lambda \quad [\mathrm{g/cm}^2],
  \label{eq:column}
\end{equation}
where $dl/d\lambda = \sqrt{g_{ij}(dx^i/d\lambda)(dx^j/d\lambda)}$ is the proper spatial arc element in the rest frame of the torus fluid (approximated here in the coordinate basis).

\subsection{Key Parameters}

\begin{center}
\begin{tabular}{lll}
\toprule
Parameter & Default & Description \\
\midrule
$a$ (spin) & 0.8 & Kerr spin, $0 \le a < 1$ \\
$r_\mathrm{obs}$ & 80\,$r_g$ & Observer radial position \\
$\theta_\mathrm{obs}$ & 70° & Observer polar angle (inclination) \\
FOV & 60° & Camera full field of view \\
$N_x \times N_y$ & 8$\times$8 & Pixel resolution \\
$r_\mathrm{max}$ & 120\,$r_g$ & Ray integration outer boundary \\
$h$ & 0.05\,$r_g$ & Integration step size \\
\bottomrule
\end{tabular}
\end{center}

\textbf{Output:} Binary file \texttt{kerr\_geodesics\_e2e.bin} in KGEO format (magic \texttt{0x4B47454F}), containing one \texttt{PathPoint} record per geodesic step: $(r,\theta,\phi,\Lambda,\tau_\mathrm{GBW},\tau_\mathrm{IIM})$.

% ─────────────────────────────────────────────────────────────────────────────
\section{Step 2: Sampling Interaction Points on Geodesics}
\label{sec:interaction_points}

\subsection{Physical Role}

The UHE neutrino travels along a null geodesic toward the black hole and has a probability of interacting with the torus material. This step samples $N_\mathrm{events}$ interaction points from the set of traced geodesics, using the accumulated column density and the DIS cross section to compute the interaction probability at each point.

\subsection{Optical Depth and Interaction Probability}

The optical depth to interaction at affine parameter $\lambda$ along a geodesic is:
\begin{equation}
  \tau(E_\nu, \lambda) = \sigma_\mathrm{CC}(E_\nu)\,\Lambda(\lambda),
  \label{eq:tau}
\end{equation}
where $\sigma_\mathrm{CC}(E_\nu)$ is the neutrino--nucleon charged-current cross section. The probability that the neutrino first interacts between $\lambda$ and $\lambda + d\lambda$ is:
\begin{equation}
  P_\mathrm{int}(\lambda)\,d\lambda = e^{-\tau(\lambda)}\,\sigma_\mathrm{CC}\,\rho(\lambda)\,dl\,d\lambda.
  \label{eq:pint}
\end{equation}
The cumulative interaction probability after traversing the full column $\Lambda$ is:
\begin{equation}
  P_\mathrm{total} = 1 - e^{-\sigma_\mathrm{CC}\,\Lambda}.
\end{equation}

Points are sampled stochastically using a fixed seed for reproducibility; each selected point carries its BL coordinates $(r,\theta,\phi)$, column $\Lambda$, and $\tau$.

\subsection{Torus Density Profile}

The torus density field $\rho(r,\theta)$ driving Eq.~\eqref{eq:column}--\eqref{eq:tau} is implemented in \texttt{src/torus\_profile.cpp}. The default is a Gaussian torus:
\begin{equation}
  \rho(r,\theta) = \rho_0 \exp\!\left(-\frac{(r - r_0)^2}{2\sigma_r^2}\right)
                  \exp\!\left(-\frac{(\theta - \pi/2)^2}{2(H/r)^2}\right),
  \label{eq:gaussian_torus}
\end{equation}
with $r_0$ the torus radius, $\sigma_r$ the radial width, and $H/r$ the aspect ratio. Additional profiles include power-law and collapsar NDAF (neutrino-dominated accretion flow) \cite{Popham1999, DiMatteo2002}.

\textbf{Output:} \texttt{interaction\_points.jsonl} --- each record contains \texttt{event\_id, x\_rg, y\_rg, z\_rg, r\_rg, theta\_rad, phi\_rad, column\_g\_cm2, tau\_gbw, tau\_iim, P\_int\_gbw, P\_int\_iim}.

% ─────────────────────────────────────────────────────────────────────────────
\section{Step 3: NLO DIS Event Generation (POWHEG + PYTHIA 8)}
\label{sec:dis}

\subsection{Physical Role}

At each sampled interaction point, a UHE neutrino undergoes charged-current deep inelastic scattering with a nucleon in the torus:
\begin{equation}
  \nu_\mu + N \;\longrightarrow\; \mu^- + X,
  \label{eq:dis}
\end{equation}
where $X$ is the hadronic final state. The cross section is calculated at NLO in QCD using POWHEG-DIS \cite{Alioli2010, POWHEG_DIS}, and the parton-level events are then showered and hadronized with PYTHIA~8 \cite{Sjostrand2014}.

\subsection{DIS Kinematics}

Define the standard DIS variables:
\begin{align}
  q^\mu &= k^\mu - k'^\mu \quad \text{(4-momentum transfer)},\\
  Q^2   &= -q^2 > 0 \quad \text{(virtuality)},\\
  x     &= \frac{Q^2}{2\,P\cdot q} \in (0,1] \quad \text{(Bjorken-}x\text{)},\\
  y     &= \frac{P\cdot q}{P\cdot k} \quad \text{(inelasticity)},\\
  W^2   &= (P + q)^2 = M_N^2 + Q^2\frac{1-x}{x} \quad \text{(invariant mass of }X\text{)},
\end{align}
where $k, k'$ are the initial and final lepton momenta, $P$ is the nucleon momentum, and $M_N$ is the nucleon mass. The center-of-mass energy squared is:
\begin{equation}
  s = 2\,M_N\,E_\nu \quad (M_N \ll E_\nu).
\end{equation}

\subsection{DIS Cross Sections: GBW and IIM Saturation Models}

At UHE energies, the small-$x$ regime dominates and parton saturation becomes relevant. HADROS-CASCADE employs two saturation models for the neutrino--nucleon cross section.

\subsubsection{Golec-Biernat--W\"{u}sthoff (GBW) Model}

The GBW dipole model \cite{GBW1998, GBW1999} parametrizes the $\gamma^*$--proton cross section via the dipole picture:
\begin{equation}
  \sigma_{\gamma^*p}(x, Q^2) = \sigma_0 \left[1 - \exp\!\left(-\frac{Q^2}{4\,Q_s^2(x)}\right)\right],
  \label{eq:gbw}
\end{equation}
where the saturation scale is:
\begin{equation}
  Q_s^2(x) = Q_0^2 \left(\frac{x_0}{x}\right)^\lambda.
\end{equation}
Fitted parameters: $\sigma_0 = 23.03$ mb, $\lambda = 0.288$, $x_0 = 3.04\times10^{-4}$, $Q_0 = 1$ GeV.

\subsubsection{Iancu--Itakura--McLerran (IIM) Model}

The IIM model \cite{IIM2002} introduces anomalous dimension corrections to the dipole amplitude:
\begin{equation}
  \mathcal{N}(r, x) =
  \begin{cases}
    \mathcal{N}_0 \left(\frac{r^2\,Q_s^2}{4}\right)^{\gamma_s + \frac{\ln(2/r\,Q_s)}{\kappa\,\lambda\,Y}} & r\,Q_s \le 2 \\
    1 - e^{-a\ln^2(b\,r\,Q_s)} & r\,Q_s > 2
  \end{cases}
  \label{eq:iim}
\end{equation}
where $Y = \ln(1/x)$, $\gamma_s = 0.6275$, $\kappa = 9.9$, and $\mathcal{N}_0 = 0.558$.

The total $\nu$--$N$ cross section at NLO is obtained by convolving these dipole cross sections with the CC structure functions and numerically integrating over $x$ and $Q^2$ up to $Q_\mathrm{max}^2 = \min(2\sqrt{M_N E_\nu},\,10^5\,\mathrm{GeV})^2$.

\subsection{POWHEG-DIS Event Generation}

The POWHEG method \cite{Nason2004, Frixione2007} generates events with NLO accuracy and a positive-weight Sudakov form factor. The channel used is charged current (CC, type 3). LHAPDF set \texttt{NNPDF31\_nlo\_as\_0118} (ID 303400) \cite{NNPDF2017} provides the parton distribution functions.

The POWHEG-DIS output is a Les~Houches Event (LHE) file with parton-level 4-momenta. These are passed to PYTHIA~8 via \texttt{Beams:frameType~=~4} (LHE input mode).

\subsection{PYTHIA 8 Hadronization}

PYTHIA~8 \cite{Sjostrand2014} applies:
\begin{enumerate}
  \item \textbf{Initial-state radiation (ISR):} backward parton shower from the struck quark.
  \item \textbf{Final-state radiation (FSR):} forward shower of the outgoing parton cascade.
  \item \textbf{Hadronization:} Lund string fragmentation model \cite{Andersson1983}, producing a realistic multiparticle hadronic final state.
  \item \textbf{Unstable hadron decays:} $\pi^0 \to \gamma\gamma$, $K_S^0 \to \pi^+\pi^-$, etc.
\end{enumerate}

\textbf{Output:} \texttt{hadros\_particle\_events.jsonl} --- one record per final-state particle with fields \texttt{event\_id, pdg, energy\_gev, px\_gev, py\_gev, pz\_gev, mass\_gev, interaction\_type}.

% ─────────────────────────────────────────────────────────────────────────────
\section{Step 4: Hadronic Cascade in GEANT4}
\label{sec:geant4}

\subsection{Physical Role}

The particles produced by PYTHIA~8 are injected into a local Cartesian box centered at the interaction point in the torus. GEANT4 \cite{Agostinelli2003, Allison2006} transports each particle, tracking electromagnetic and hadronic secondary production, until particles either escape the box or deposit all their energy.

\subsection{Particle Classification}

Before sending particles to GEANT4, the pipeline classifies each final-state particle from PYTHIA into categories:

\begin{center}
\begin{tabular}{lll}
\toprule
Category & Criterion & Fate \\
\midrule
\texttt{invisible} & PDG $\in \{12, 14, 16, -12, -14, -16\}$ & Passes through; counted as $E_\mathrm{invisible}$ \\
\texttt{untracked} & Zero momentum or unknown PDG & Skipped; counted in $E_\mathrm{untracked}$ \\
\texttt{unsupported\_uhe} & $E_k > E_k^\mathrm{max}$ (per species) & Flagged; counted as $E_\mathrm{unsupported}$ \\
\texttt{transportable} & Everything else & Transported through GEANT4 \\
\bottomrule
\end{tabular}
\end{center}

The \texttt{unsupported\_uhe} category is the most common for UHE primaries: GEANT4 hadronic physics lists (FTFP\_BERT, QGSP\_BERT) have upper energy limits of $\sim10^5$~GeV, above which the nuclear interactions are not modeled. These particles are flagged as \texttt{unsupported\_uhe\_energy} in the energy budget. This is why the GEANT4 step typically returns a \texttt{PARTIAL} status for UHE runs.

\subsection{ZAMO Local Tetrad Frame}

The local GEANT4 box is defined in the Locally Non-Rotating Frame (LNRF), also called the Zero-Angular-Momentum Observer (ZAMO) frame \cite{Bardeen1972}. The orthonormal basis vectors at position $(r_\mathrm{int}, \theta_\mathrm{int}, \phi_\mathrm{int})$ are:
\begin{equation}
  \hat{e}_{\hat{r}} = \frac{1}{\sqrt{g_{rr}}}\,\partial_r, \qquad
  \hat{e}_{\hat{\theta}} = \frac{1}{\sqrt{g_{\theta\theta}}}\,\partial_\theta, \qquad
  \hat{e}_{\hat{\phi}} = \frac{1}{\sqrt{g_{\phi\phi}}}\,\partial_\phi,
  \label{eq:tetrad}
\end{equation}
with the Kerr metric components evaluated at the interaction point:
\begin{equation}
  \sqrt{g_{rr}} = \sqrt{\Sigma/\Delta}, \qquad
  \sqrt{g_{\theta\theta}} = \sqrt{\Sigma}, \qquad
  \sqrt{g_{\phi\phi}} = \sqrt{A}\sin\theta/\sqrt{\Sigma}.
\end{equation}

Momenta of PYTHIA particles are rotated into this local frame before being passed to GEANT4. After the cascade, escaped particles are transformed back to BL coordinates using the inverse tetrad.

\subsection{Energy Conservation}

The energy budget at the interaction point is:
\begin{equation}
  E_\mathrm{in} = E_\mathrm{dep} + E_\mathrm{esc} + E_\mathrm{invisible} + E_\mathrm{unsupported} + E_\mathrm{untracked} + E_\mathrm{failed} + E_\mathrm{residual},
  \label{eq:energy_budget}
\end{equation}
where:
\begin{itemize}
  \item $E_\mathrm{dep}$: energy deposited in the box (heating torus material),
  \item $E_\mathrm{esc}$: energy carried by escaped particles (observable),
  \item $E_\mathrm{invisible}$: neutrinos passing through (from DIS final state),
  \item $E_\mathrm{unsupported}$: UHE particles above GEANT4 energy limit,
  \item $E_\mathrm{untracked}$: particles with unknown or zero momentum,
  \item $E_\mathrm{failed}$: GEANT4 jobs that errored,
  \item $E_\mathrm{residual}$: floating-point closure error.
\end{itemize}

The run achieves \texttt{VALIDATED} status if $|E_\mathrm{residual}| < 1$~GeV and all jobs complete; otherwise it is \texttt{PARTIAL} (warning, not error).

\subsection{Physics Lists}

GEANT4 is configured with either:
\begin{itemize}
  \item \textbf{FTFP\_BERT:} Fritiof Precompound model + Bertini cascade. Standard for $E_k \gtrsim 5$~GeV.
  \item \textbf{QGSP\_BERT:} Quark--Gluon String Precompound + Bertini below 12~GeV.
\end{itemize}

\textbf{Output:} \texttt{geant4\_ready\_particles.jsonl} --- each escaped particle record contains \texttt{pdg, energy\_gev, global\_exit\_x\_rg, global\_exit\_y\_rg, global\_exit\_z\_rg}, plus the BL position and 4-momentum in the Kerr frame.

% ─────────────────────────────────────────────────────────────────────────────
\section{Step 5: Particle-Ray Association Camera}
\label{sec:camera}

\subsection{Physical Role}

Each particle that escapes the GEANT4 box now has a global position in Kerr spacetime. The particle-ray association camera records which sampled Kerr ray is spatially close to that particle and, in the angular mode, locally aligned with its momentum. This is a cascade-origin / particle-ray association map, not a detector image and not full transport of secondary particles to the observer.

\subsection{Backward Ray Tracing}

For each pixel $(i,j)$, a null geodesic $\gamma_{ij}(\lambda)$ is integrated backward from the observer. The particle at position $\mathbf{x}_p = (r_p, \theta_p, \phi_p)$ is associated with pixel $(i,j)$ if:
\begin{equation}
  \min_\lambda \left[ d_\mathrm{spatial}(\gamma_{ij}(\lambda),\,\mathbf{x}_p) \right] < \epsilon_s
  \quad \text{and} \quad
  d_\mathrm{angular}(\hat{n}_{ij},\,\hat{n}_p) < \epsilon_a,
  \label{eq:matching}
\end{equation}
where $\epsilon_s$ and $\epsilon_a$ are spatial and angular matching tolerances (configured in \texttt{config\_web\_final.py}), and $\hat{n}_p$ is the local particle momentum direction at the association point. In \texttt{spatial\_only} mode only the spatial criterion is applied and the direction misalignment is recorded as not calculated.

The associated energy in pixel $(i,j)$ is redshift-adjusted for the sampled Kerr ray:
\begin{equation}
  E_\mathrm{obs} = E_p \cdot \frac{(p_\mu u^\mu)_\mathrm{source}}{(p_\mu u^\mu)_\mathrm{observer}},
  \label{eq:redshift}
\end{equation}
where $u^\mu$ is the 4-velocity of the emitter (ZAMO frame) and observer (static at $r_\mathrm{obs}$).

\textbf{Output:} \texttt{particle\_ray\_association\_camera.csv}. Legacy \texttt{observed\_particles\_by\_pixel.*} files may also be written for compatibility; the legacy name does not imply a physical observation.

% ─────────────────────────────────────────────────────────────────────────────
\section{Step 6: Ray--Event Linking}
\label{sec:ray_link}

\subsection{Physical Role}

After the association step, each ray-associated secondary particle can be linked to the specific Kerr geodesic sample through which the incident neutrino traveled. This ``ray link'' allows propagating the incoming column density $\Lambda_\mathrm{in}$ from the pre-computed geodesic cache to every downstream record.

\subsection{Modular Assignment}

The pipeline generates $N_\mathrm{POWHEG}$ events (typically 16) but may have only $N_\mathrm{int}$ sampled interaction points (e.g.\ 10). Event IDs are assigned modularly:
\begin{equation}
  \text{link}(e) =
  \begin{cases}
    \text{links}[e] & \text{if } e \in \text{links} \\
    \text{links}[\text{sorted\_keys}[(e-1)\,\mathrm{mod}\,|\text{links}|]] & \text{otherwise}
  \end{cases}
  \label{eq:modular}
\end{equation}

This ensures every POWHEG event gets a valid geodesic, even when the number of events exceeds the number of sampled interaction points.

\textbf{Output:} Updated JSONL files with \texttt{incoming\_ray\_id}, \texttt{column\_before\_g\_cm2}, \texttt{tau\_before\_gbw}, \texttt{tau\_before\_iim} fields added to each record.

% ─────────────────────────────────────────────────────────────────────────────
\section{Step 7: GBW/IIM Interaction Probability Weighting}
\label{sec:weighting}

\subsection{Physical Role}

Not all geodesics have the same probability of producing a DIS interaction. This step reweights each ray-associated secondary-particle row by the interaction probability derived from the real incoming geodesic column density:
\begin{equation}
  w_\mathrm{GBW}(e) = P_\mathrm{int}^\mathrm{GBW}(e) = 1 - \exp\!\left(-\sigma_\mathrm{GBW}(E_\nu)\,\Lambda_\mathrm{in}(e)\right),
  \label{eq:weight_gbw}
\end{equation}
\begin{equation}
  w_\mathrm{IIM}(e) = P_\mathrm{int}^\mathrm{IIM}(e) = 1 - \exp\!\left(-\sigma_\mathrm{IIM}(E_\nu)\,\Lambda_\mathrm{in}(e)\right).
  \label{eq:weight_iim}
\end{equation}

Cross sections $\sigma_\mathrm{GBW}(E_\nu)$ and $\sigma_\mathrm{IIM}(E_\nu)$ are read from precomputed tables in \texttt{sigma\_table\_path}. The weighted image is:
\begin{equation}
  \mathcal{I}_{ij} = \sum_{e: (i_e, j_e) = (i,j)} E_\mathrm{obs}(e)\,w(e),
  \label{eq:image}
\end{equation}
which represents a weighted cascade-origin / particle-ray association map. It is not a detector flux image and does not transport secondary particles to the observer.

\textbf{Output:} \texttt{gbw\_iim\_camera\_summary.csv} with per-event weights and \texttt{weighted\_energy\_gev} column.

% ─────────────────────────────────────────────────────────────────────────────
\section{Step 8: Scientific Figures and Dashboard}
\label{sec:output}

\subsection{Generated Products}

The final steps produce:

\begin{enumerate}
  \item \textbf{Particle-ray association RGB map} (\texttt{particle\_ray\_association\_rgb.png}): false-color association map with gamma (red), electromagnetic (green), hadronic (blue) channels.
  \item \textbf{Per-channel energy maps}: separate images for $\gamma$, $e^\pm$, hadrons, $\nu$.
  \item \textbf{Energy budget chart}: pie/bar plot showing $E_\mathrm{dep}$, $E_\mathrm{esc}$, $E_\mathrm{invisible}$, etc.
  \item \textbf{GBW vs.\ IIM sensitivity}: side-by-side comparison of the two saturation models.
  \item \textbf{Camera sensitivity scan}: varied $a$, $\theta_\mathrm{obs}$, $r_\mathrm{obs}$, FOV, step size.
  \item \textbf{Interactive dashboard} (\texttt{dashboard/index.html}): HTML page aggregating all run products with clickable plots.
\end{enumerate}

\subsection{Geodesic Camera Preview}

A separate, non-science tool (\texttt{hadros\_geodesic\_preview}) renders a low-resolution Kerr geodesic image using only the geometric optics (no DIS, no cascade). It supports interactive navigation via GLFW/OpenGL and saves camera configurations to \texttt{configs/cameras/last\_camera.json}. Supported display modes: \texttt{celestial\_sphere}, \texttt{disk\_intersection}, \texttt{shadow\_only}, \texttt{combined}.

% ─────────────────────────────────────────────────────────────────────────────
\section{Data Flow Summary}
\label{sec:dataflow}

\begin{longtable}{p{2.2cm}p{3.5cm}p{3.5cm}p{4cm}}
\toprule
\textbf{Step} & \textbf{Input} & \textbf{Output} & \textbf{Format} \\
\midrule
\endhead
1 Kerr Geodesics & Metric params, camera config & \texttt{kerr\_geodesics\_e2e.bin} & Binary KGEO (PathPoint stream) \\
\midrule
2 Interaction Points & KGEO binary & \texttt{interaction\_points.jsonl} & JSONL \\
\midrule
3 POWHEG+PYTHIA & \texttt{interaction\_points.jsonl} + POWHEG binary & \texttt{hadros\_particle\_events.jsonl} & JSONL + CSV \\
\midrule
4 GEANT4 Cascade & \texttt{hadros\_particle\_events.jsonl} & \texttt{geant4\_ready\_particles.jsonl} & JSONL (with global positions) \\
\midrule
5 Backward Camera & \texttt{geant4\_ready\_particles.jsonl} & \texttt{observed\_particles\_by\_pixel.csv} & CSV \\
\midrule
6 Ray--Event Link & All above + KGEO & Updated JSONL files with \texttt{incoming\_ray\_id} & JSONL \\
\midrule
7 GBW/IIM Weight & Linked JSONL + $\sigma$ tables & \texttt{gbw\_iim\_camera\_summary.csv} & CSV \\
\midrule
8 Figures & All CSVs/JSONLs & PNG images + HTML dashboard & PNG, HTML \\
\bottomrule
\caption{HADROS-CASCADE pipeline data flow.}
\label{tab:dataflow}
\end{longtable}

% ─────────────────────────────────────────────────────────────────────────────
\section{Configuration and Execution}
\label{sec:config}

\subsection{Configuration File}

The pipeline is controlled by a single JSON file (default: \texttt{presets/config\_web/final\_config.json}). Key sections:

\begin{lstlisting}[language={}]
{
  "run":              { "run_name", "output_dir", "physics_mode" },
  "black_hole_camera":{ "spin", "black_hole_mass_msun",
                        "observer_inclination_deg", "field_of_view_deg",
                        "resolution", "observer_radius_rg",
                        "ray_max_radius_rg", "ray_step" },
  "uhe_dis":          { "source_model", "energy_gev", "dis_model",
                        "spectral_bins", "sigma_table",
                        "powheg_executable", "pythia8_config",
                        "n_events", "seed" },
  "torus":            { "rho0", "r0_rg", "sigma_rg", "h_over_r",
                        "density_profile" },
  "outputs":          { "science_plots", "dashboard",
                        "validation_plots", "diagnostic_plots" }
}
\end{lstlisting}

\subsection{Key Make Targets}

\begin{center}
\begin{tabular}{ll}
\toprule
Target & Action \\
\midrule
\texttt{make build} & Compile all C++ binaries \\
\texttt{make kerr-geodesics} & Run Step~1 \\
\texttt{make final-config-web} & Start configuration web UI \\
\texttt{make final-pipeline-dry-run} & Print planned steps \\
\texttt{make final-pipeline-run} & Execute full pipeline \\
\texttt{make geodesic\_preview} & Launch real Kerr geodesic preview \\
\texttt{make cascade\_geant4\_local\_box} & Compile GEANT4 binary \\
\bottomrule
\end{tabular}
\end{center}

\subsection{Resumability}

Each GEANT4 job (one particle per job) is independently resumable. If a job fails or the run is interrupted, completed jobs are read from the \texttt{jobs/} directory and skipped on the next run, so only failed or missing jobs are re-executed. This is controlled by \texttt{run\_powheg\_pythia\_geant4\_resumable.py}.

% ─────────────────────────────────────────────────────────────────────────────
\section{Exit Codes and Status Flags}
\label{sec:status}

\begin{center}
\begin{tabular}{lll}
\toprule
Exit Code & Status String & Meaning \\
\midrule
0 & \texttt{VALIDATED} & All criteria met; physics reliable \\
2 & \texttt{PARTIAL} & Outputs present but $<$ validation threshold; \\
  &                   & typical for UHE runs dominated by \texttt{unsupported\_uhe} \\
$\neq 0,2$ & \texttt{FAILED} & Hard error; pipeline stops \\
\bottomrule
\end{tabular}
\end{center}

The \texttt{[warn] geant4\_real\_safe\_zamo: partial return code 2 with required outputs present} message means the GEANT4 step exited with code 2 (\texttt{PARTIAL}), but all required output files exist and the pipeline continues. This is the expected behavior for UHE neutrino energies ($E_\nu \gtrsim 10^5$~GeV) where most DIS secondary particles exceed the GEANT4 hadronic physics range.

% ─────────────────────────────────────────────────────────────────────────────
\section*{Abbreviations}

\begin{tabular}{ll}
BL & Boyer--Lindquist \\
CC & Charged Current \\
DIS & Deep Inelastic Scattering \\
FOV & Field of View \\
GBW & Golec-Biernat--W\"{u}sthoff \\
GEANT4 & Geometry And Tracking (CERN simulation toolkit) \\
IIM & Iancu--Itakura--McLerran \\
JSONL & JSON Lines (one JSON object per line) \\
LHE & Les Houches Event format \\
LNRF & Locally Non-Rotating Frame \\
NLO & Next-to-Leading Order \\
POWHEG & Positive-Weight Hardest Emission Generator \\
PYTHIA & Parton-level event generator \\
QCD & Quantum Chromodynamics \\
RK4/RK5 & Runge--Kutta 4th/5th order \\
UHE & Ultra-High-Energy \\
ZAMO & Zero-Angular-Momentum Observer \\
\end{tabular}

% ─────────────────────────────────────────────────────────────────────────────
\begin{thebibliography}{99}

\bibitem{Kerr1963}
R.~P.~Kerr, \textit{Gravitational Field of a Spinning Mass as an Example of Algebraically Special Metrics}, Phys.\ Rev.\ Lett.\ \textbf{11}, 237 (1963).

\bibitem{Misner1973}
C.~W.~Misner, K.~S.~Thorne, J.~A.~Wheeler, \textit{Gravitation}, W.H.~Freeman, San Francisco (1973).

\bibitem{MTW}
C.~W.~Misner, K.~S.~Thorne, J.~A.~Wheeler, \textit{Gravitation}, W.H.~Freeman, San Francisco (1973).

\bibitem{Carter1968}
B.~Carter, \textit{Global Structure of the Kerr Family of Gravitational Fields}, Phys.\ Rev.\ \textbf{174}, 1559 (1968).

\bibitem{Bardeen1972}
J.~M.~Bardeen, W.~H.~Press, S.~A.~Teukolsky, \textit{Rotating Black Holes: Locally Nonrotating Frames, Energy Extraction, and Scalar Synchrotron Radiation}, Astrophys.\ J.\ \textbf{178}, 347 (1972).

\bibitem{Fuerst2004}
S.~V.~Fuerst, K.~Wu, \textit{Radiation transfer of emission lines in curved spacetime}, Astron.\ Astrophys.\ \textbf{424}, 733 (2004).

\bibitem{GBW1998}
K.~Golec-Biernat, M.~W\"{u}sthoff, \textit{Saturation effects in deep inelastic scattering at low $x$ and its implications on diffraction}, Phys.\ Rev.\ D \textbf{59}, 014017 (1998). [arXiv:hep-ph/9807401]

\bibitem{GBW1999}
K.~Golec-Biernat, M.~W\"{u}sthoff, \textit{Saturation in diffractive deep inelastic scattering}, Phys.\ Rev.\ D \textbf{60}, 114023 (1999). [arXiv:hep-ph/9903358]

\bibitem{IIM2002}
E.~Iancu, K.~Itakura, S.~McLerran, \textit{Geometric scaling above the saturation scale}, Nucl.\ Phys.\ A \textbf{708}, 327 (2002). [arXiv:hep-ph/0203137]

\bibitem{Alioli2010}
S.~Alioli, P.~Nason, C.~Oleari, E.~Re, \textit{A general framework for implementing NLO calculations in shower Monte Carlo programs: the POWHEG BOX}, JHEP \textbf{06}, 043 (2010). [arXiv:1002.2581]

\bibitem{POWHEG_DIS}
S.~Alioli, P.~Nason, C.~Oleari, E.~Re, \textit{NLO vector-boson production matched with shower in POWHEG}, JHEP \textbf{07}, 060 (2008); see also POWHEG-DIS implementation for $\nu N$ scattering.

\bibitem{Nason2004}
P.~Nason, \textit{A new method for combining NLO QCD with shower Monte Carlo algorithms}, JHEP \textbf{11}, 040 (2004). [arXiv:hep-ph/0409146]

\bibitem{Frixione2007}
S.~Frixione, P.~Nason, C.~Oleari, \textit{Matching NLO QCD computations with Parton Shower simulations: the POWHEG method}, JHEP \textbf{11}, 070 (2007). [arXiv:0709.2092]

\bibitem{Sjostrand2014}
T.~Sj\"{o}strand et al., \textit{An Introduction to PYTHIA~8.2}, Comput.\ Phys.\ Commun.\ \textbf{191}, 159 (2015). [arXiv:1410.3012]

\bibitem{Andersson1983}
B.~Andersson, G.~Gustafson, G.~Ingelman, T.~Sj\"{o}strand, \textit{Parton Fragmentation and String Dynamics}, Phys.\ Rept.\ \textbf{97}, 31 (1983).

\bibitem{NNPDF2017}
NNPDF Collaboration, \textit{Parton distributions from high-precision collider data}, Eur.\ Phys.\ J.\ C \textbf{77}, 663 (2017). [arXiv:1706.00428]

\bibitem{Agostinelli2003}
S.~Agostinelli et al.\ (GEANT4 Collaboration), \textit{GEANT4 — a simulation toolkit}, Nucl.\ Instrum.\ Meth.\ A \textbf{506}, 250 (2003).

\bibitem{Allison2006}
J.~Allison et al., \textit{GEANT4 Developments and Applications}, IEEE Trans.\ Nucl.\ Sci.\ \textbf{53}, 270 (2006).

\bibitem{Popham1999}
R.~Popham, S.~E.~Woosley, C.~Fryer, \textit{Hyperaccreting Black Holes and Gamma-Ray Bursts}, Astrophys.\ J.\ \textbf{518}, 356 (1999). [arXiv:astro-ph/9807028]

\bibitem{DiMatteo2002}
T.~Di~Matteo, R.~Perna, R.~Narayan, \textit{Neutrino Trapping and Emission in Neutrino-dominated Accretion Flows}, Astrophys.\ J.\ \textbf{579}, 706 (2002). [arXiv:astro-ph/0207319]

\end{thebibliography}

\end{document}
